\section{Covariant baseline theory}
The baseline action is
\begin{equation}
S=\int d^4x\sqrt{-g}\left[\frac{\Mpl^2}{2}R-\frac12 g^{\mu\nu}\nabla_\mu\phi\nabla_\nu\phi-V(\phi)\right]+S_m+S_r.
\end{equation}
A periodic stratification potential is parameterized as
\begin{equation}
V(\phi)=V_{\rm off}+V_0\left[1-\cos\left(\frac{N_s\phi}{f}\right)\right].
\end{equation}
Here $[\phi]=[f]=\mathrm{mass}$, $[V_0]=[V_{\rm off}]=\mathrm{mass}^4$, and $N_s$ is a positive integer. Variation with respect to the metric gives
\begin{align}
\Mpl^2G_{\mu\nu}&=T^{(m)}_{\mu\nu}+T^{(r)}_{\mu\nu}+T^{(\phi)}_{\mu\nu},\\
T^{(\phi)}_{\mu\nu}&=\nabla_\mu\phi\nabla_\nu\phi-g_{\mu\nu}\left[\frac12(\nabla\phi)^2+V(\phi)\right],
\end{align}
and variation with respect to $\phi$ gives
\begin{equation}
\Box\phi-V_{,\phi}=0,\qquad
V_{,\phi}=\frac{N_sV_0}{f}\sin\left(\frac{N_s\phi}{f}\right).
\end{equation}
Minimal coupling and the matter equations imply $\nabla_\mu T^{\mu\nu}_{(i)}=0$ for each noninteracting component. This conservation law is a consequence of diffeomorphism invariance rather than an optional numerical closure.

For any null vector $n^\mu$, the canonical scalar obeys
\begin{equation}
T^{(\phi)}_{\mu\nu}n^\mu n^\nu=(n^\mu\nabla_\mu\phi)^2\ge0.
\end{equation}
Consequently the baseline cannot use scalar null-energy-condition violation to obtain a spatially flat regular bounce. The model has a controlled GR limit and no higher-derivative degrees of freedom, but those virtues also make its bounce mechanism highly constrained.

None of $V_{\rm off}$, $V_0$, or $f$ is derived by the action itself. $V_{\rm off}$ is a renormalized vacuum-energy parameter; changing it to produce a desired cycle is a tuning, not a derivation. $N_s$ can be discrete by definition, but its value is not fixed by four-dimensional covariance. Computing these quantities from more fundamental data would require an additional microscopic symmetry or ultraviolet completion.
